\setcounter{figure}{12}
\begin{figure}[t!]
\centering
\includegraphics[width=\textwidth]{fig_spatial_maps.pdf}
\caption{\textbf{Envelope and quadratic spatial maps.} Illustrative line profile:
$\varepsilon_1(x)=\varepsilon_0(1+x/L_x)$ and
$\varepsilon_2(x)=\varepsilon_0(1-x/L_x)$ on $-L_x\le x\le L_x$, with polarization scale
$\varepsilon_0>0$ (mV) and length $L_x>0$ (m). In (a,b), each geometric curve is
divided by its own positive value at $x=0$.
\emph{(a)} Modulation is $1-|x/L_x|$ for envelope following and $1-(x/L_x)^2$ for quadratic
detection: both peak centrally, but the quadratic profile is broader here.
\emph{(b)} The mean envelope and mean-square polarization $\mathcal I$ both
increase toward the edges; neither includes a signed detector gain.
\emph{(c)} The static JR rate source $Y_{\rm DC}=\Sig''\mathcal I/2$ for uniform
positive curvature (solid red) or negative curvature of the same magnitude
(dashed blue). Both use the same positive scale $|\Sig''|\varepsilon_0^2/2$,
giving $\pm[1+(x/L_x)^2]$; the gray line marks zero. These are assumed detector
states before recurrent feedback, not inferred physiological responses.
This is a line profile, not a head model or a demonstration of three-dimensional focusing.}
\label{fig:spatial_maps}
\end{figure}
